\section*{Hoofdstuk \ref{ch:b2:alkene-redox} --- Oxidatie en reductie van alkenen}

\begin{solution}{exo:b2:alkene-redox:1}
Methylcyclohexaan; 1-methylcyclohexaan-1-ol; trans-2-methylcyclohexaan-1-ol (syn-additie van
H en B, dus de \ce{OH} en de nieuwe H aan dezelfde kant, de methylgroep trans ten opzichte van de \ce{OH});
1-methylcyclohexeenoxide; cis-1-methylcyclohexaan-1,2-diol; 6-oxoheptanal
\ce{CH3CO(CH2)4CHO}.
\end{solution}

\begin{solution}{exo:b2:alkene-redox:2}
Hydroborering--oxidatie: 2-methylpropaan-1-ol. Zuurgekatalyseerde hydratatie:
2-methylpropaan-2-ol.
\end{solution}

\begin{solution}{exo:b2:alkene-redox:3}
trans-2,3-Dimethyloxiraan, de twee methylgroepen aan tegenovergestelde kanten van de ring. Het is chiraal (geen
spiegelvlak) en wordt als racemisch mengsel gevormd.
\end{solution}

\begin{solution}{exo:b2:alkene-redox:4}
Reductieve opwerking: twee moleculen propanal. Oxidatieve opwerking: twee moleculen
propaanzuur.
\end{solution}

\begin{solution}{exo:b2:alkene-redox:5}
(a) meso-butaan-2,3-diol, inactief omdat het achiraal is. (b) de
(2\textit{R},3\textit{R})- en (2\textit{S},3\textit{S})-diolen in gelijke hoeveelheden, inactief
omdat ze racemisch zijn.
\end{solution}

\begin{solution}{exo:b2:alkene-redox:6}
In zuur valt water het tertiaire koolstofatoom aan, in base valt hydroxide de \ce{CH2} aan; in beide
gevallen komt er een \ce{OH} op elk koolstofatoom: hetzelfde 2-methylpropaan-1,2-diol. Met methanol geven de twee
routes verschillende ethers (zoals in het hoofdstuk).
\end{solution}

\begin{solution}{exo:b2:alkene-redox:7}
Hydrogenering met één equivalent \ce{H2} over een vergiftigde palladiumkatalysator
(Lindlar): de syn-additie aan de drievoudige binding stopt bij het (\textit{Z})-alkeen.
\end{solution}

\begin{solution}{exo:b2:alkene-redox:8}
Propanon (3 C) en butanal (4 C) geven \ce{(CH3)2C=CH-CH2CH2CH3},
2-methylhex-2-een: ze worden verbonden via hun carbonylkoolstofatomen.
\end{solution}

\begin{solution}{exo:b2:alkene-redox:9}
\ce{BH3} (of een omvangrijk boraan), daarna \ce{H2O2} en \ce{NaOH}: hexaan-1-ol. Zuurgekatalyseerde hydratatie
verloopt via het secundaire carbokation en geeft hexaan-2-ol (Markovnikov), met
omgelegde bijproducten.
\end{solution}

\begin{solution}{exo:b2:alkene-redox:10}
Het \omterm{def:b2:alkene-redox:epoxide}{peroxyzuur}, een elektrofiel, reageert sneller met de meer gesubstitueerde, elektronenrijkere
dubbele binding in de ring. Het omvangrijke boraan reageert met de minder gehinderde eindstandige \ce{C=CH2}.
\end{solution}

\begin{solution}{exo:b2:alkene-redox:11}
trans: \omterm{def:b2:alkene-redox:epoxide}{peroxyzuur}, daarna water in zuur (anti-opening). cis: katalytisch \ce{OsO4} met een
heroxidator, of koud verdund permanganaat (syn).
\end{solution}

\begin{solution}{exo:b2:alkene-redox:12}
\ce{C6H10} heeft onverzadigingsgraad 2; één equivalent \ce{H2} betekent één \ce{C=C}, dus
één ring. Eén enkel dialdehyde met zes koolstofatomen komt van een ring die bij zijn dubbele binding geopend is:
cyclohexeen.
\end{solution}

\begin{solution}{pb:b2:alkene-redox:1}
\textbf{1.} $10 \times 12.011 + 16 \times 1.008 = \qty{136.24}{g/mol}$.
\textbf{2.} $(2 \times 10 + 2 - 16)/2 = 3$.
\textbf{3.} Eén ring en twee dubbele bindingen.
\textbf{4.} \ce{C10H16 + 2H2 -> C10H20}, 1-isopropyl-4-methylcyclohexaan.
\textbf{5.} $1.00/136.24 = \qty{7.34}{mmol}$ limoneen, \qty{14.68}{mmol} \ce{H2}.
\textbf{6.} $V = nRT/p = 0.01468 \times 8.314 \times 298.15/10^5 = \qty{3.64e-4}{m^3}$,
\qty{364}{mL}.
\textbf{7.} Nee: de ringkoolstofatomen 1 en 4 dragen elk twee identieke ringpaden; het cis- en het
trans-isomeer zijn achiraal.
\textbf{8.} Methanal \ce{HCHO} (1 C) uit het uiteinde \ce{=CH2}, en één enkele \ce{C9}-verbinding,
3-acetyl-6-oxoheptanal \ce{CH3CO-CH2CH2-CH(COCH3)-CH2-CHO}.
\textbf{9.} Haar twee koolstofatomen behoren tot dezelfde ring: de dubbele binding verbreken opent de ring
zonder iets af te splitsen.
\textbf{10.} Methanal (een gas).
\textbf{11.} De aldehydegroepen worden carbonzuren (methanal gaat verder tot \ce{CO2}); de
ketonen blijven onveranderd.
\textbf{12.} Het keton en het aldehyde van de \ce{C9}-keten waren in de ring verbonden; het tweede
keton (\ce{COCH3}) en methanal waren de twee uiteinden van de dubbele binding van de zijketen.
\textbf{13.} Eén (C4 van de ring); de \omterm{def:b2:alkene-redox:cleavage}{ozonolyse} raakt het niet, en het blijft over als
stereocentrum van het \ce{C9}-product.
\textbf{14.} De exocyclische \ce{C=CH2}, minder gehinderd dan de drievoudig gesubstitueerde dubbele binding in de ring.
\textbf{15.} 2-(4-Methylcyclohex-3-een-1-yl)propaan-1-ol: \ce{CH2OH} op het vroegere eindstandige
koolstofatoom.
\textbf{16.} Primair.
\textbf{17.} De \ce{OH} gaat naar het minst gesubstitueerde koolstofatoom, tegengesteld aan de oriëntatie
die de regel van Markovnikov voor hydratatie voorspelt.
\textbf{18.} Twee configuraties aan het nieuwe centrum, wat twee diastereomeren geeft (met het bestaande
centrum).
\textbf{19.} Zuurgekatalyseerde hydratatie (of een andere markovnikovhydratatie): de tertiaire alcohol,
$\alpha$-terpineol.
\textbf{20.} De dubbele binding in de ring: drievoudig gesubstitueerd, elektronenrijker.
\textbf{21.} 1,2-Epoxy-1-methyl-4-(prop-1-een-2-yl)cyclohexaan, met het zuurstofatoom als brug tussen C1 en C2.
\textbf{22.} Op C1, het meest gesubstitueerde koolstofatoom (dat de methylgroep draagt).
\textbf{23.} 1-Methyl-4-(prop-1-een-2-yl)cyclohexaan-1,2-diol, de twee \ce{OH} trans (anti-opening).
\textbf{24.} \omterm{def:b2:alkene-redox:dihydroxylation}{Syn-dihydroxylering} van de dubbele binding in de ring (katalytisch \ce{OsO4}), wat het cis-diol
geeft --- als het reagens selectief met die dubbele binding kan reageren.
\textbf{25.} \qty{1.00}{g} limoneen neemt $\boldsymbol{\approx \qty{364}{mL}}$
diwaterstof op.
\end{solution}
